\chapter{Collisions in 3D}

\section{The same idea, one more axis}

The 3D equivalent of a \lstinline|Rectangle| is an \textbf{AABB}: an
\emph{axis-aligned bounding box}. A box that never rotates, described by its two
opposite corners.

\begin{lstlisting}
typedef struct BoundingBox {
    Vector3 min;    // the corner with the smaller x, y and z
    Vector3 max;    // the corner with the larger  x, y and z
} BoundingBox;
\end{lstlisting}

And the rule is chapter 6's rule with a third line added: two boxes collide if
they overlap on \textbf{all three} axes at once.

\begin{lstlisting}
overlapX = (a.min.x <= b.max.x) && (a.max.x >= b.min.x);
overlapY = (a.min.y <= b.max.y) && (a.max.y >= b.min.y);
overlapZ = (a.min.z <= b.max.z) && (a.max.z >= b.min.z);
hit      = overlapX && overlapY && overlapZ;
\end{lstlisting}

raylib gives you it as \lstinline|CheckCollisionBoxes(a, b)|. Its relatives:

\begin{lstlisting}
CheckCollisionBoxes(boxA, boxB);
CheckCollisionSpheres(centreA, rA, centreB, rB);
CheckCollisionBoxSphere(box, centre, radius);
\end{lstlisting}

\section{Centre and size, not corners}

Nobody stores their objects as two corners. You store a position and a size, and
convert when you need to test:

\begin{lstlisting}
BoundingBox BoxFromCenter(Vector3 centre, Vector3 size)
{
    BoundingBox b;
    b.min = (Vector3){ centre.x - size.x/2, centre.y - size.y/2, centre.z - size.z/2 };
    b.max = (Vector3){ centre.x + size.x/2, centre.y + size.y/2, centre.z + size.z/2 };
    return b;
}
\end{lstlisting}

You will write this function in every 3D project you ever start. It is the same
``subtract half the size'' from chapter 3, with a z on the end.

\gotcha{The player's \lstinline|camera.position| is the height of their
\textbf{eyes}, not their middle. A 1.7 m person standing on the floor has eyes at
$Y = 1.7$ and a body centred at $Y = 0.85$. Build the collision box around the
body, not the eyes, or your player will collide with things at head height and
walk through everything below the knee.}

\begin{lstlisting}
Vector3 bodyCentre = { position.x, position.y - EYE_HEIGHT/2.0f, position.z };
Vector3 bodySize   = { RADIUS*2, EYE_HEIGHT, RADIUS*2 };
\end{lstlisting}

\section{Resolving, per axis, exactly as in 2D}

\begin{lstlisting}
Vector3 attempt = camera.position;
attempt.x += dx;
if (Collides(attempt)) { /* refused: keep the old x */ }
else camera.position.x = attempt.x;

attempt = camera.position;
attempt.z += dz;
if (Collides(attempt)) { /* refused: keep the old z */ }
else camera.position.z = attempt.z;
\end{lstlisting}

X then Z, because on foot those are the two ground axes. If you add jumping and
gravity, Y becomes a third pass with the same shape --- and the Y pass is where
you discover you have landed, because a refused downward move means there is
floor under you.

\keyidea{Detect with one call, resolve one axis at a time. It was true in 2D, it
is true in 3D, and it will still be true in whatever engine you use next.}

\section{Debugging: draw the boxes}

\begin{lstlisting}
DrawBoundingBox(playerBox, GREEN);
DrawBoundingBox(wallBox, hit ? RED : GREEN);
\end{lstlisting}

Nine times out of ten, a collision bug is not a logic bug --- the box is simply
not where you think it is. It is offset by half a height, or it is built from
last frame's position, or it is twice the size of the thing you drew. Draw it and
the bug is visible in a second instead of an hour.

\section{Grids: the other way, and the fast one}

Testing against a list of boxes is fine for a few dozen walls. For a whole level
there is a better approach, and it is what the original Doom did: store the level
as a \textbf{grid}, and convert a position straight into a cell index.

\begin{lstlisting}
static bool IsWall(float x, float z)
{
    int col = (int)floorf(x/CELL + 0.5f);
    int row = (int)floorf(z/CELL + 0.5f);

    if (col < 0 || col >= MAP_COLS || row < 0 || row >= MAP_ROWS) return true;

    return (MAP[row][col] == '#');
}
\end{lstlisting}

No loop. It costs the same whether the map has 10 walls or 10,000, and the
out-of-bounds check doubles as an invisible wall around the level.

To give the player some width, test the four corners of the square around them
instead of the single centre point:

\begin{lstlisting}
static bool CanStandAt(float x, float z)
{
    const float r = PLAYER_RADIUS;
    return !IsWall(x - r, z - r) && !IsWall(x + r, z - r)
        && !IsWall(x - r, z + r) && !IsWall(x + r, z + r);
}
\end{lstlisting}

\section{Shooting: rays}

A \lstinline|Ray| is a starting point and a direction:

\begin{lstlisting}
typedef struct Ray { Vector3 position; Vector3 direction; } Ray;
\end{lstlisting}

Fire one out of the camera and ask what it meets:

\begin{lstlisting}
Ray shot = { camera.position, forward };
RayCollision hit = GetRayCollisionBox(shot, EnemyBox(enemy.position));

if (hit.hit) { /* hit.distance, hit.point and hit.normal are filled in */ }
\end{lstlisting}

For a mouse click rather than a crosshair, build the ray from the cursor:

\begin{lstlisting}
Ray pick = GetScreenToWorldRay(GetMousePosition(), camera);
\end{lstlisting}

\gotcha{A ray goes on forever and hits everything in its path. To find what you
actually shot, loop over the candidates, keep the one with the smallest
\lstinline|hit.distance|, and compare it against the distance to the nearest
wall. Skip that last comparison and your players will happily shoot enemies
through solid walls.}

\tryit{Lesson 5 shows one box turning red on contact, with the three axis
overlaps printed separately; lesson 6 is a closed room where the walls actually
stop you.}{./course3d/build.sh 5 \quad and \quad ./course3d/build.sh 6}
